% Данный файл распространяется под лицензией CC BY 4.0. Текст лицензии размещён на https://creativecommons.org/licenses/by/4.0/
% (c) Кафедра системного программирования, 2026

% Компилировать XeTeX

\documentclass[a4paper]{article}

\usepackage[a4paper, top=8mm, bottom=8mm, left=8mm, right=8mm]{geometry}

\usepackage{polyglossia}
\setdefaultlanguage[babelshorthands=true]{russian}
\setotherlanguage{english}

\usepackage{fontspec}
\setmainfont{FreeSerif}
\newfontfamily{\russianfonttt}[Scale=0.7]{DejaVuSansMono}

\usepackage[tiny, compact]{titlesec}

\usepackage{titling}
\setlength{\droptitle}{-1cm}
\pretitle{\begin{center}\begin{bfseries}\Large}
\posttitle{\par\end{bfseries}\end{center}}
\preauthor{\begin{center}\normalsize}
\postauthor{\par\end{center}\vspace{-1.8cm}}

\usepackage{hyperref}
\usepackage{bookmark}
\usepackage{csquotes}

\title{Виртуальное блочное устройство на базе оперативной памяти в ядре Linux}

\author{Климович Станислав Андреевич}

\date{}

\begin{document}

\maketitle

\begin{flushright}
    Группа: \emph{25.Б41-мм}

    Кафедра: \emph{системного программирования}

    Научный руководитель: \emph{Васенина Анна Игоревна}

    Номер семестра практики: \emph{2}
\end{flushright}

\section{Введение}
В Unix-подобных операционных системах всё является файлом:
блочное устройство представляет собой накопитель данных или его раздел, доступный через виртуальную файловую систему (VFS) и пригодный
для монтирования с файловой системой поверх. Такая абстракция позволяет
единообразно работать с реальным оборудованием и программными (виртуальными)
устройствами. Виртуальные блочные устройства применяются в
облачных средах, системах виртуализации и встроенных системах~--- там, где
реального дискового оборудования нет или оно избыточно. Понимание их
внутреннего устройства необходимо разработчикам storage-стека, гипервизоров
и авторам вкладов в ядро Linux.

В рамках данной учебной практики реализован модуль ядра \texttt{/dev/ram\_virtblk}~---
простейший RAM-блочник (базовый уровень сложности в иерархии виртуальных
блочных устройств курса)~--- с CI/CD-инфраструктурой и автоматизированным
интеграционным тестированием на реальном ядре. Полученные навыки применимы
при разработке драйверов, storage-подсистем и участии в open-source проектах
Linux-экосистемы.

\section{Постановка задачи}
\textbf{Целью работы является} разработка модуля ядра Linux, реализующего
виртуальное блочное устройство \texttt{/dev/ram\_virtblk}, поддерживающее
форматирование, монтирование и файловые операции с хранением данных в RAM.

\textbf{Для достижения цели требуется решить следующие задачи:}
\begin{enumerate}
    \item Изучить подсистему блочных устройств ядра Linux и существующие
          реализации RAM-дисков.
    \item Реализовать модуль ядра на основе BIO-based API с настроенной
          CI/CD-инфраструктурой: линтингом и автоматической генерацией документации.
    \item Верифицировать корректность устройства интеграционными тестами
          на целевом ядре (\texttt{6.18.13-200.fc43.x86\_64}).
\end{enumerate}

\section{Обзор}

Цель обзора~--- выбрать архитектурный подход для реализации и оценить
применимость существующих решений в качестве основы.

В ядре Linux существуют штатные RAM-диски: \texttt{brd} хранит данные
в radix-tree из страниц памяти и поддерживает разделы; \texttt{zram} дополнительно
сжимает данные алгоритмами lz4/zstd; \texttt{loop}-устройство монтирует файлы
образов как блочные устройства. Все используют многоочередной стек blk-mq и
представляют собой производственный код объёмом от нескольких сотен до тысяч строк.

Архитектура блочного устройства в ядре предполагает два подхода к обработке I/O:
классический на основе очереди запросов (blk-mq) и упрощённый BIO-based --- прямой
обработчик \texttt{submit\_bio} без промежуточной диспетчеризации.
BIO-based API появился в Linux~5.9 и значительно проще для изучения; для
RAM-устройства он не создаёт потерь производительности, так как отсутствует
аппаратная латентность, ради которой blk-mq и разрабатывалось.

\textbf{Вывод:} \texttt{brd} и \texttt{zram} избыточны для учебной задачи и
скрывают именно те механизмы, которые необходимо освоить. BIO-based API выбран
как оптимальный: минимальный по сложности и достаточный по функциональности.

\section{Описание предлагаемого решения}

\textbf{Структура модуля.} Три файла: \texttt{virtblk.h}~--- константы устройства
(имя, размер сектора, количество секторов, суммарный объём 64~МиБ); \texttt{virtblk.c}~---
реализация; \texttt{Makefile}~--- сборка через стандартный механизм Kbuild.

\textbf{Управление памятью.} RAM-буфер (131072~сектора $\times$ 512~байт) выделяется
через \texttt{vzalloc()}. От \texttt{kmalloc()} отказались: 64~МиБ не требуют
физически непрерывной памяти; \texttt{vzalloc()} работает с виртуально
непрерывными, но физически разрозненными страницами, снижая давление на аллокатор ядра.

\textbf{Поток I/O.} Ядро вызывает \texttt{ram\_virtblk\_submit\_bio()}
с BIO-запросом. Функция вычисляет байтовое смещение (\texttt{bi\_sector}~$\times$~512),
проверяет выход за границы (при превышении~--- \texttt{bio\_io\_error()}), затем
обходит сегменты через \texttt{bio\_for\_each\_segment()}. Для каждого сегмента
\texttt{ram\_virtblk\_transfer()} маппит физическую страницу через
\texttt{kmap\_local\_page()} (thread-safe замена устаревшего \texttt{kmap()}),
выполняет \texttt{memcpy()} в/из буфера и немедленно размаппирует через
\texttt{kunmap\_local()}, предотвращая утечки маппингов.

\textbf{Инициализация и выгрузка} реализованы через цепочку \texttt{goto}:
\texttt{vzalloc()} $\to$ \texttt{register\_blkdev()} $\to$ \texttt{blk\_alloc\_disk()}
$\to$ \texttt{add\_disk()}. При ошибке на шаге $N$ все предыдущие шаги
откатываются в обратном порядке. Флаг \texttt{GENHD\_FL\_NO\_PART} запрещает
сканирование таблицы разделов~--- устройство монолитно, без разделов.

\textbf{CI/CD.} Два воркфлоу GitHub Actions. \texttt{lint.yml}: \texttt{checkpatch.pl}~---
официальный линтер ядра Linux~--- проверяет стиль кода; Doxygen
генерирует документацию из исходных комментариев и публикует на GitHub Pages.
\texttt{test.yml}: сначала проверяется точная версия ядра на runner-е, затем
собирается и запускается \texttt{.ko}.

\section{Тестирование}

Модульное тестирование неприменимо к коду ядра: модуль исполняется в kernel~space
и не запускается в изоляции. KASAN (аналог ASAN для ядра) требует пересборки
с \texttt{CONFIG\_KASAN=y}, что выходит за рамки задачи. Применены четыре
интеграционных теста на реальном устройстве (Fedora~Server~43, ядро
\texttt{6.18.13-200.fc43.x86\_64}, гипервизор VirtualBox):

\begin{enumerate}
    \item \textbf{Базовый write/read}~--- запись строки через \texttt{dd}
          в сырое устройство, чтение и побайтовое сравнение.
    \item \textbf{Файловая система}~--- форматирование в ext4 (\texttt{mkfs.ext4}),
          монтирование, запись и чтение файла через VFS.
    \item \textbf{Целостность данных}~--- запись блока случайных данных из
          \texttt{/dev/urandom}, побайтовое сравнение снапшота через \texttt{cmp}.
    \item \textbf{Граница устройства}~--- I/O-запрос за конец устройства;
          ожидается вызов \texttt{bio\_io\_error()} и ошибка на стороне ядра.
\end{enumerate}

Тесты запускаются автоматически при каждом пуше через GitHub Actions и проходят
успешно: устройство корректно выполняет I/O, монтируется как ext4, сохраняет
побайтовую целостность данных и возвращает ошибку при выходе за границы.

\section{Заключение}

В ходе выполнения работы получены следующие результаты:

\begin{itemize}
    \item Проведён обзор существующих реализаций RAM-дисков ядра Linux (brd, zram,
          loop) и выбран BIO-based подход как оптимальный для учебной задачи.
    \item Реализован модуль ядра Linux (\texttt{virtblk.c}, $\approx$150~строк)
          на основе BIO-based API: регистрация блочного устройства, обработчик
          \texttt{submit\_bio}, инициализация и выгрузка с обработкой ошибок.
    \item Зарегистрировано виртуальное блочное устройство \texttt{/dev/ram\_virtblk}
          ёмкостью 64~МиБ, поддерживающее форматирование в ext4, монтирование и
          стандартные файловые операции.
    \item Настроен CI/CD-пайплайн: линтинг \texttt{checkpatch.pl} и автогенерация
          документации Doxygen на GitHub-hosted runner; интеграционные тесты на
          self-hosted runner с целевой версией ядра.
    \item Написаны и автоматизированы четыре интеграционных теста, покрывающих
          базовый I/O, файловую систему, целостность данных и граничные случаи.
\end{itemize}

Репозиторий: \url{https://github.com/Monrealle/linux-virtblk}

Техническая документация: \url{https://monrealle.github.io/linux-virtblk/}

\begin{thebibliography}{9}

\bibitem{lovyagin}
Ловягин~Н.~Ю., Ловягин~Ю.~Н. Основы работы в Unix-подобных операционных системах:
учебное пособие. Ч.~I.~--- СПб., 2020.~--- 157~с.~---
(§1.8 Файловая система; §E.7 Монтирование разделов; §E.8 Работа с разделами и файловыми системами). ~--- URL:~\url{https://api.dspace.spbu.ru/server/api/core/bitstreams/b95862a3-b6a5-45b8-9a68-283a7291dbbf/content}.

\bibitem{ldd3}
Corbet~J., Rubini~A., Kroah-Hartman~G. Linux Device Drivers [Электронный ресурс].~---
3-е изд.~--- O'Reilly Media, 2005.~--- 636~с.~--- (Главы 1, 2, 4, 16).~---
URL:~\url{https://lwn.net/Kernel/LDD3/}.

\bibitem{codestyle}
Linux kernel coding style // The Linux Kernel documentation. — URL:~\url{https://docs.kernel.org/process/coding-style.html}

\end{thebibliography}

\end{document}
